\documentclass[11pt]{article}
\usepackage[T1]{fontenc}
\usepackage{lmodern}
\usepackage[margin=1.04in]{geometry}
\usepackage{amsmath,amssymb,amsthm,mathrsfs,booktabs,microtype}
\usepackage[colorlinks=true,linkcolor=blue,citecolor=blue,urlcolor=blue]{hyperref}
\numberwithin{equation}{section}
\newtheorem{theorem}{Theorem}[section]
\newtheorem{proposition}[theorem]{Proposition}
\newtheorem{lemma}[theorem]{Lemma}
\newtheorem{corollary}[theorem]{Corollary}
\theoremstyle{remark}\newtheorem{remark}[theorem]{Remark}
\newcommand{\Q}{\mathbb Q}\newcommand{\Z}{\mathbb Z}
\newcommand{\Qbar}{\overline{\mathbb Q}}
\newcommand{\eps}{\varepsilon}\newcommand{\LL}{\mathcal L}
\newcommand{\HH}{\mathcal H}\newcommand{\wt}{\widehat}
\DeclareMathOperator{\spanop}{span}
\title{Window bounds for repetitions and block complexity\\of algebraic numbers}
\author{}\date{}
\hypersetup{pdftitle={Window bounds for repetitions and block complexity of algebraic numbers},pdfsubject={Research manuscript: fixed-support moving weights and dyadic quality layers}}
\begin{document}\maketitle
\begin{abstract}
Let $\xi$ be a real algebraic irrational number and $b\ge2$ an integer.
Write $R(L)$ for the length of the shortest prefix of its fractional
base-$b$ expansion containing two occurrences of some word of length $L$,
and $p(L)$ for its block complexity. For every fixed $\theta>1$ there is
$c>0$ such that every sufficiently large interval $[X,X^\theta]$ contains
an integer $L$ for which
\[
 R(L)>cL\sqrt{\frac{\log L}{\log\log L}},\qquad
 p(L)>\frac c2L\sqrt{\frac{\log L}{\log\log L}}.
\]
Consequently the logarithmically weighted complexity limsup is infinite
for every weight exponent $\gamma>1/2$. We also obtain windows
$[X,X^{1+o(1)}]$ below this endpoint. The arithmetic input is a
finite-support moving-weight sparse-parameter theorem proved here by
adapting the auxiliary-polynomial argument of Evertse and Ferretti.
The additional combinatorial step retains all dyadic prefix scales:
when several scales yield the same approximation scale, their
multiplicity gives an exponential improvement in approximation quality.
Summing fixed-quality counting bounds removes a logarithmic loss without
requiring an extension with a different smallness exponent at each point.
\end{abstract}

\section{Introduction}
Fix a real algebraic irrational number $\xi$ and an integer $b\ge2$.
We may assume $0<\xi<1$, and write
$\xi=\sum_{i\ge1}d_ib^{-i}$, with $d_i\in\{0,\ldots,b-1\}$.
Define
\[
 p(L)=\#\{d_{j+1}\cdots d_{j+L}:j\ge0\}.
\]
Following Bugeaud and Kim \cite{BK19}, let $R(L)$ be the least prefix
length containing two, possibly overlapping, occurrences of a word
of length $L$. We use $R$ rather than $r$ to distinguish this function
from the preperiod length of a rational approximation. Pigeonholing gives
\begin{equation}\label{eq:RvsP}
 R(L)\le L+p(L).
\end{equation}
Indeed, the prefix of length $L+p(L)$ contains $p(L)+1$ starting
positions for length-$L$ words. All logarithms without a subscript
are natural, and iterated logarithms are used only for sufficiently
large arguments.

\begin{theorem}\label{thm:main}
For every fixed $\theta>1$, there is a constant
$c=c(\xi,b,\theta)>0$ such that, for every sufficiently large real $X$,
there is an integer $L\in[X,X^\theta]$ satisfying
\begin{equation}\label{eq:main}
 R(L)>cL\sqrt{\frac{\log L}{\log\log L}},\qquad
 p(L)>\frac c2L\sqrt{\frac{\log L}{\log\log L}}.
\end{equation}
\end{theorem}

\begin{corollary}\label{cor:limsup}
For each $g\in\{R,p\}$,
\[
 \limsup_{L\to\infty}
 \frac{g(L)\sqrt{\log\log L}}{L\sqrt{\log L}}>0.
\]
For every $\gamma>1/2$ and every $\eta<1/2$,
\begin{equation}\label{eq:weightedlimsup}
 \limsup_{L\to\infty}
 \frac{g(L)(\log\log L)^\gamma}{L\sqrt{\log L}}=\infty,
 \qquad
 \limsup_{L\to\infty}\frac{g(L)}{L(\log L)^\eta}=\infty.
\end{equation}
\end{corollary}

The repetition-function statement is proved directly. Inequality
\eqref{eq:RvsP} then transfers its lower bound to $p$; a lower bound
for $p$ alone would not imply the repetition-function statement.

\begin{theorem}\label{thm:nearwindow}
Let $F(x)=(\log x)^u(\log\log x)^{-\gamma}$, where $u>0$, $\gamma\ge0$,
and either $u<1/2$, or $u=1/2$ and $\gamma>1/2$.
For every $C>0$ there is $A=A(\xi,b,C,u,\gamma)>0$ such that, for every
sufficiently large $X$, the interval
\begin{equation}\label{eq:nearwindow}
 \left[X,\ X\exp\{A F(X)^2\log(2F(X))\}\right]
\end{equation}
contains an integer $L$ with both $R(L)>CLF(L)$ and $p(L)>CLF(L)$.
Its upper endpoint is $X^{1+o(1)}$.
\end{theorem}

Adamczewski and Bugeaud proved $p(L)/L\to\infty$ for algebraic
irrationals in integer bases \cite{AB07}. Bugeaud and Evertse proved
the logarithmic limsup bound with $\eta<1/11$
\cite[Theorem~2.1]{BE08}. The arithmetic argument here uses their
repetition method and the semistable machinery of
\cite[Sections~8--14]{EF13}. A fixed-weight theorem alone does not
permit the exact ratio $\rho=r/t$ to vary from point to point.
Section~\ref{sec:moving} proves the needed finite-support variant,
including a bound for its height threshold; Section~\ref{sec:stable}
checks semistability for every weight in the relevant family.

The subsequent counting argument uses that theorem separately at each
fixed quality. If $\mu$ dyadic prefix scales yield a single dyadic
approximation scale, the largest prefix yields quality at least a
constant times $2^\mu/B$. Summation over the multiplicity levels gives
a bound $O(B^2\log(2B))$ for the original number of scales.
This replaces greedy thinning of the prefix scales.
The proof of this step is in Section~\ref{sec:layers}.

Constants may depend on $\xi,b$, and on the fixed window or function
parameters explicitly indicated. No effective numerical starting
length or assertion uniform in varying algebraic numbers or bases
is made. The endpoint statement is a positive limsup; it does not
assert a positive lower bound at every sufficiently large length.

\section{From repetitions to approximations}\label{sec:words}
\begin{lemma}\label{lem:words}
Suppose that the prefix of length $H$ contains two occurrences of a
word of length $k$. Then there are integers $r,s,t,v,a$ such that
\begin{equation}\label{eq:repeat}
 t=r+s,\quad 0\le r<t,\quad t+v\le H,\quad s\ge v\ge k/3,
 \qquad 0\le a\le b^t-b^r,
\end{equation}
and
\begin{equation}\label{eq:approx}
 |(b^t-b^r)\xi-a|\le b^{-v}.
\end{equation}
\end{lemma}
\begin{proof}
Let $D$ be the difference of the two starting positions. If $D\ge k$,
the two occurrences are disjoint and we take $v=k$. If $D<k$, their
union is a word of length $k+D$ with period $D$. Its first two blocks
of length $v'=D\lfloor(k+D)/(2D)\rfloor$ are equal and disjoint.
Here $v'\ge k/3$: if $D\ge k/3$, use $v'\ge D$; otherwise use
$v'>(k-D)/2$. Thus in either case the prefix has a decomposition
$UVWVX$, with $v=|V|\ge k/3$.

Put $r=|U|$, $s=|VW|$, $t=r+s$, and
$A_j=\sum_{i=1}^jd_ib^{j-i}$, $A_0=0$. Set $a=A_t-A_r$.
The number $a/(b^t-b^r)$ has expansion $U(VW)^\infty$ and agrees with
$\xi$ through the first $t+v$ digits. Its distance from $\xi$ is at most
$b^{-t-v}$, which proves \eqref{eq:approx}. If $P$ is the integer value
of $VW$, then $a=(b^s-1)A_r+P$, giving $0\le a\le b^t-b^r$.
The closed interval of values with a specified finite prefix also
covers a periodic tail consisting entirely of the digit $b-1$.
The remaining length bounds follow from the disjoint occurrences.
\end{proof}

\begin{lemma}[Repetitions in a prescribed window]\label{lem:windowwords}
Let $2\le X<Y$, let $F$ be positive and nondecreasing on $[X,Y]$, and
write $F_Y=F(Y)$. Suppose
\[
 R(k)\le CkF(k)\quad\text{for every integer }k\in[X,Y],
 \qquad CF_Y\ge1.
\]
For each integer $i$ with $8CF_YX\le2^i\le Y$, there is an
approximation satisfying \eqref{eq:repeat}--\eqref{eq:approx}, with
$H=2^i$, such that
\begin{equation}\label{eq:windowwords}
 v_i\ge\frac{2^i}{24CF_Y},\qquad t_i\ge X/3.
\end{equation}
\end{lemma}
\begin{proof}
Take $k_i=\lfloor2^i/(4CF_Y)\rfloor$. The lower bound on $2^i$ gives
$k_i\ge2X-1\ge X$, and $CF_Y\ge1$ gives $k_i\le Y$.
Thus $R(k_i)\le CF_Yk_i\le2^i/4$, so the prefix of length $2^i$
contains the required repetition. Also
$k_i\ge2^i/(8CF_Y)$. Lemma~\ref{lem:words} gives
$v_i\ge k_i/3\ge2^i/(24CF_Y)$. Since $t_i\ge v_i$,
the bound $t_i\ge X/3$ follows. All uses of the assumed upper bound
for $R$ occur at lengths inside $[X,Y]$.
\end{proof}

\section{Twisted heights and the exact weights}\label{sec:height}
For a number field $K$ write $M_K$ for its places, normalized by
$\|z\|_v=|z|_p^{d_v}$ on $z\in\Q$ when $v\mid p$, where
$d_v=[K_v:\Q_p]/[K:\Q]$ and $\sum_{v\mid p}d_v=1$.
At an infinite place represented by an embedding $\sigma_v$ the
corresponding formula is $\|z\|_v=|\sigma_v(z)|^{d_v}$.
For a system of linearly independent forms
$\LL=(L_{iv}:v\in M_K,1\le i\le3)$ and real weights $c_{iv}$, put
\begin{equation}\label{eq:height}
 H_{\LL,c,Q}(x)=\prod_{v\in M_K}
 \max_{1\le i\le3}\bigl(\|L_{iv}(x)\|_vQ^{-c_{iv}}\bigr).
\end{equation}
For $x\in\Qbar^3$ use a finite extension $E/K$ containing its coordinates,
keep the forms unchanged at $w\mid v$, and multiply the weights by
$d(w\mid v)=[E_w:K_v]/[E:K]$. This gives an extension-independent height.

Choose a finite Galois extension $K/\Q$ in $\mathbb C$ containing
$\xi$, and let $G=\operatorname{Gal}(K/\Q)$ and
$d=[\Q(\xi):\Q]\ge2$. Write $\beta_1,\ldots,\beta_d$ for the distinct
conjugates of $\xi$. They are all nonzero.
For each infinite place $v$, choose $\sigma_v\in G$ with
$\|z\|_v=|\sigma_v(z)|^{d_v}$, and set $\beta_v=\sigma_v^{-1}(\xi)$.
Define the fixed forms by
\begin{equation}\label{eq:forms}
 (L_{1v},L_{2v},L_{3v})=
 \begin{cases}(X_1,X_2,L_{\beta_v}),&v\mid\infty,\\
 (X_1,X_2,X_3),&v\nmid\infty,
 \end{cases}
 \quad L_\beta=-\beta X_1+\beta X_2+X_3.
\end{equation}
Their local determinants are all $1$, and there are $d+3$ distinct forms.
For every rational vector $x$,
\begin{equation}\label{eq:transport}
 \|L_{\beta_v}(x)\|_v
 =|\sigma_v(L_{\beta_v}(x))|^{d_v}=|L_\xi(x)|^{d_v}.
\end{equation}
This is the inverse-conjugate transport used in
\cite[Section~5, equations (5.14)--(5.17)]{BE08}.
It uses smallness only at the specified real value of $\xi$.

\begin{lemma}[Mass of a conjugate]\label{lem:mass}
For every conjugate $\beta$ of $\xi$,
\begin{equation}\label{eq:mass}
 \sum_{\substack{v\mid\infty\\\beta_v=\beta}}d_v=\frac1d.
\end{equation}
\end{lemma}
\begin{proof}
Let $\iota\in G$ be complex conjugation on $K$ and
$H=\langle\iota\rangle$. Infinite places are represented by the right
cosets $H\sigma$. Since $\xi$ is real,
$(\iota\sigma)^{-1}(\xi)=\sigma^{-1}(\xi)$, so $\beta_v$ does not depend
on the representative. The field $K$ is either totally real or totally
imaginary, and $d_v=|H|/|G|$ in either case.
The set of $\sigma\in G$ with $\sigma^{-1}(\xi)=\beta$ has cardinality
$|G|/d$ by orbit--stabilizer. It is a union of right $H$-cosets, each of
cardinality $|H|$. Its corresponding total place weight is therefore
$(|G|/(d|H|))(|H|/|G|)=1/d$.
\end{proof}

For each rational prime $p\mid b$, put
\begin{equation}\label{eq:primeweights}
 \omega_p=\frac{\operatorname{ord}_p(b)\log p}{\log b},\qquad
 \kappa_v=d_v\omega_p\quad(v\mid p).
\end{equation}
These numbers are positive and satisfy
$\sum_{p\mid b}\omega_p=1$ and
$\sum_{v\mid b}\kappa_v=1$, where $v\mid b$ means $v\mid p$ for some
$p\mid b$. Let $S$ consist of all infinite places and all places above
the prime divisors of $b$.

Fix $0<\eps\le1/2$ and $0\le\rho\le1-\eps$. In the order of
\eqref{eq:forms}, define raw weights by
\begin{equation}\label{eq:raw}
 e_v=\begin{cases}
 d_v(1,\rho,-\eps),&v\mid\infty,\\
 \kappa_v(-1,-\rho,0),&v\mid b,\\
 (0,0,0),&v\notin S.
 \end{cases}
\end{equation}
Put
\begin{equation}\label{eq:center}
 \overline e_v=\tfrac13\sum_i e_{iv},\quad D_\eps=1+\eps/3,
 \quad c_{iv}(\rho)=\frac{e_{iv}-\overline e_v}{D_\eps},
 \quad \delta=\frac{\eps}{3+\eps}.
\end{equation}
The two mass identities give
\begin{equation}\label{eq:norm}
 \sum_i c_{iv}=0,\qquad \sum_v\max_i c_{iv}=1,
 \qquad \sum_v\overline e_v=-\eps/3.
\end{equation}
Indeed $\sum_{v,i}e_{iv}=-\eps$ and the sum of the raw local maxima
is $1$; centering changes the latter to $1+\eps/3$.

Consider any approximation satisfying \eqref{eq:repeat}--\eqref{eq:approx}.
Choose $0<\eps\le\min(v/t,1/2)$ and put $\rho=r/t$.
Since $t-r\ge v$, we have $\rho\le1-\eps$.
Let $x=(b^t,b^r,a)$, $\Psi=b^t$, and $Q=\Psi^{D_\eps}$.
At the infinite places, \eqref{eq:approx} and \eqref{eq:transport} give
$\|L_{iv}(x)\|_v\le\Psi^{e_{iv}(\rho)}$.
At a place $v\mid p\mid b$, the same bounds follow from
\[
 \|b^t\|_v=\Psi^{-\kappa_v},\qquad
 \|b^r\|_v=\Psi^{-\kappa_v\rho},\qquad \|a\|_v\le1.
\]
At all other finite places the coordinates are integral. Hence
\begin{equation}\label{eq:small}
 H_{\LL,c(\rho),Q}(x)
 \le\Psi^{\sum_v\overline e_v}=Q^{-\delta},
 \qquad \delta=\frac{\eps}{3+\eps}.
\end{equation}
For a collection of such approximations with a common $\eps$, the
condition $t_{j+1}>2t_j$ implies $Q_{j+1}>Q_j^2$.
No primitivity of $x$ is needed: $\Psi$ is an external parameter,
not an identification with its projective height.

\section{Uniform semistability}\label{sec:stable}
For a $q$-dimensional subspace $U\subset\Qbar^3$, define its local
weight as the minimum of $\sum_{i\in I}c_{iv}$ over all $q$-element
sets $I$ for which $(L_{iv}|_U)_{i\in I}$ is linearly independent.
Its global weight $w_c(U)$ is the sum over places. The system is
semistable if $w_c(U)\le0$ for every nonzero proper $U$.
Define $w_e$ analogously. Centering gives
\begin{equation}\label{eq:weightshift}
 w_c(U)=\frac{w_e(U)+q\eps/3}{D_\eps}.
\end{equation}

\begin{proposition}\label{prop:stable}
For the family \eqref{eq:raw}--\eqref{eq:center}, all nonzero proper
subspaces have negative weight. More precisely,
\[
 w_c(U)\le-\frac{\eps}{6D_\eps}\quad(\dim U=1),\qquad
 w_c(U)\le-\frac{\eps}{3D_\eps}\quad(\dim U=2).
\]
\end{proposition}
\begin{proof}
By Lemma~\ref{lem:mass}, the infinite-place contribution to $w_e(U)$
can be computed using each distinct $L_\beta$ once, with raw weights
$(1,\rho,-\eps)/d$. The finite contributions combine into one minimum
calculation with coordinate forms and weights $(-1,-\rho,0)$, since
$\sum_{v\mid b}\kappa_v=1$. This grouping concerns only the subspace
weight calculation; all original places and norms remain in the height.

Consider a line $U$. If no $L_\beta$ vanishes on it, its infinite
contribution is $-\eps$ and its finite contribution is at most zero.
If two distinct $L_\beta$ vanish, subtracting them shows
$U=\langle(1,1,0)\rangle$; all $L_\beta$ then vanish and
$w_e(U)=\rho-1\le-\eps$.
Otherwise precisely one $L_\beta$ vanishes. If $X_1|_U\ne0$, the
exceptional infinite contribution is at most $1/d$, the remaining
infinite contribution is $-(1-1/d)\eps$, and the finite contribution
is $-1$. Thus $w_e(U)\le-(1-1/d)(1+\eps)$.
If $X_1|_U=0$, then $X_2|_U\ne0$ and
$w_e(U)=-(1-1/d)(\rho+\eps)$.
Since $d\ge2$, every line therefore has $w_e(U)\le-\eps/2$.

For planes write $U=\ker M$, and put
$A_\beta=\spanop(L_\beta,X_2)$ and $B=\spanop(X_1,X_2)$ in the dual
space. For distinct conjugates $\beta,\beta'$,
\[
 A_\beta\cap A_{\beta'}=A_\beta\cap B=\langle X_2\rangle.
\]
This follows by comparing the $X_1$ and $X_3$ coefficients.
At a grouped infinite place the minimum pair has weight
$(\rho-\eps)/d$ unless $M\in A_\beta$; the finite minimum is
$-1-\rho$ unless $M\in B$.
The following list is exhaustive. Proportional normals are identified.
\begin{center}\small
\begin{tabular}{@{}ll@{}}\toprule
Position of $M$ & $w_e(\ker M)$\\\midrule
$M\notin B\cup\bigcup_\beta A_\beta$ & $-1-\eps$\\
$M\in A_\beta\setminus(\langle L_\beta\rangle\cup\langle X_2\rangle)$
 & $-1-\eps+(1-\rho)/d$\\
$M\in\langle L_\beta\rangle$ for some $\beta$ & $-(1-1/d)(1+\eps)$\\
$M\in B\setminus(\langle X_1\rangle\cup\langle X_2\rangle)$ & $\rho-1-\eps$\\
$M\in\langle X_1\rangle\cup\langle X_2\rangle$ & $-\eps$\\\bottomrule
\end{tabular}
\end{center}
Each entry is at most $-\eps$, since $d\ge2$, $\eps\le1/2$ and
$\rho\le1-\eps$.
For example, at $M=L_\beta$ the pair for that conjugate must be
$(X_1,X_2)$, whereas at $M=X_1$ the finite pair is $(X_2,X_3)$.
These observations also cover the ties when $\rho=0$.
Equation \eqref{eq:weightshift} proves both bounds.
\end{proof}

\section{A finite-support moving-weight bound}\label{sec:moving}
We prove the arithmetic input separately. Throughout this section the
number field $K$, the finite set $S$, the forms $\LL$, and the finite
place $v_0\notin S$ are fixed. The dimension is three. We assume that
$S$ contains every infinite place, that the forms outside $S$ are the
coordinate forms, and that their determinant at every place is $1$.
Write $s=|S|$. We call a system of real weights admissible if
\begin{equation}\label{eq:admissible}
 c_{iv}=0\ (v\notin S),\qquad
 \sum_i c_{iv}=0,\qquad \sum_v\max_i c_{iv}\le1,
 \qquad w_c(U)\le0\ (0\ne U\subsetneq\Qbar^3).
\end{equation}
The last condition is the semistability established in
Section~\ref{sec:stable} for the weights used in our application.

\begin{theorem}\label{thm:sparse}
There are constants depending only on the fixed data with the following
property. For $0<\delta\le1$ there exist $\Omega(\delta)>2$ and
$C_*(\delta)>e$ with
\begin{equation}\label{eq:cutofforder}
 \Omega(\delta)\ll\delta^{-5},\qquad
 \log\log C_*(\delta)\ll\delta^{-2}\log(2/\delta),
\end{equation}
such that any sequence $(c^{(h)},Q_h)$ of admissible weights and parameters
with
\[
 Q_1\ge C_*(\delta),\qquad Q_{h+1}>Q_h^{\Omega(\delta)},\qquad
 \{x\ne0:H_{\LL,c^{(h)},Q_h}(x)\le Q_h^{-\delta}\}\ne\varnothing
\]
has length $O(\delta^{-2})$.
\end{theorem}

Here is the structure of the proof. At each parameter a small twisted
height produces a gap between successive infima. Semistability turns
that gap into a height lower bound for an infimum subspace. An exterior
power then gives two independent vectors with small local coordinates
in a plane of large height. After selecting a common exterior power,
we construct one auxiliary polynomial for all the selected parameters.
A nonvanishing proposition supplies a derivative that is nonzero at
grid points in the planes. Its local bounds give a product less than
$1$, contradicting the product formula.

We use the fixed-form height comparison and absolute Minkowski theorem
from \cite[Lemma~7.1 and Proposition~9.2]{EF13}, the determinant estimate
from \cite[Lemma~10.2]{EF13}, and the rescaling and Davenport construction
from \cite[Lemmas~11.1--11.3]{EF13}. For the polynomial we use the
number-field Siegel lemma \cite[Lemma~13.1]{EF13}; the nonvanishing input
is \cite[Proposition~12.1]{EF13}. Below we give the determinant
calculation, the simultaneous coefficient conditions with changing
weights, the derivative estimates, and the cutoff calculation explicitly.

All constants are uniform in $c^{(h)},Q_h,\delta$ once the fixed data
have been chosen. The exterior forms will depend only on $\LL$ and
the exterior degree. The planes, vectors, and local exponents may
depend on $h$. In particular the permutation arising from Davenport's
construction changes the assignment of exponents, not the underlying
forms. This distinction allows the coefficient equations to be imposed
over $K$ in fixed coordinate systems. We do not apply the final
fixed-weight interval theorem of \cite{EF13} separately at each point.

\subsection{Pointwise data and the height of the exterior plane}
We first give a pointwise version with a cutoff of order
$\exp(O(1/\delta))$. This smaller cutoff will be used before imposing
the height condition in the nonvanishing proposition.

\begin{lemma}[Pointwise exterior data]\label{lem:packet}
There are constants $A_0>0$ and $c_0>0$, depending only on $K,\LL$, such
that if $c$ is admissible, $Q\ge\exp(A_0/\delta)$ and
$H_{\LL,c,Q}(x)\le Q^{-\delta}$ for some nonzero $x$, then there is
$k\in\{1,2\}$ with the following properties. Put $N=\binom3k=3$.
There are fixed exterior forms $\wt L_{lv}$, depending only on $k,\LL$,
and independent vectors $z_1,z_2\in\Qbar^3$ spanning a plane $T$ such that
\begin{equation}\label{eq:packetheight}
 H_2(T)\ge Q^{c_0\delta}.
\end{equation}
For every place $w$ of a common finite extension $E/K$,
\begin{equation}\label{eq:packetlocal}
 \|\wt L_{lw}(z_j)\|_w\le
 \begin{cases}Q^{d(w\mid v)a_{lv}},&w\mid v\ne v_0,\\
 Q^{d(w\mid v_0)b_l},&w\mid v_0,
 \end{cases}
\end{equation}
for $l=1,2,3$ and $j=1,2$, where, with $\gamma_v=\max_i c_{iv}$,
\begin{equation}\label{eq:packetweights}
 \sum_l a_{lv}=0,\quad |a_{lv}|\le2\gamma_v,\quad
 a_{lv}=0\ (v\notin S),\quad
 \sum_l b_l\le-\delta/3,\quad |b_l|\le3.
\end{equation}
Here $H_2$ is the absolute Euclidean height of the exterior coordinate
vector defining a subspace, as in \cite[Section~6]{EF13}.
\end{lemma}
\begin{proof}
\emph{The gap between successive infima.}
Let $\lambda_1\le\lambda_2\le\lambda_3$ be the successive infima of
$H_{\LL,c,Q}$ over $\Qbar$. Since $\sum_v\gamma_v\le1$, the fixed-form
height comparison gives a constant $C_0\ge1$ such that
\begin{equation}\label{eq:minima}
 C_0^{-1}Q^{-1}\le\lambda_1,\qquad
 a_*:=3^{-3/2}\le P_\lambda:=\lambda_1\lambda_2\lambda_3\le8.
\end{equation}
The second assertion is the absolute Minkowski theorem; the determinant
assumption makes its determinant factor equal to $1$.
The assumed small vector gives $\lambda_1\le Q^{-\delta}$. Hence
$a_*\le Q^{-\delta}\lambda_3^2$, so
\[
 \lambda_3\ge a_*^{1/2}Q^{\delta/2}\ge1
 \quad\text{provided }Q^\delta\ge3^{3/2}.
\]
Consequently
$(\lambda_1/\lambda_2)(\lambda_2/\lambda_3)\le Q^{-\delta}$.
For some $k\in\{1,2\}$ we therefore have
\begin{equation}\label{eq:minimagap}
 g:=\lambda_k/\lambda_{k+1}\le Q^{-\delta/2}<1.
\end{equation}
By \cite[Lemma~9.1]{EF13}, the points below any height strictly between
$\lambda_k$ and $\lambda_{k+1}$ span the same $k$-dimensional subspace
$T_k$, defined over $K$.

Put $A_k=\lambda_1\cdots\lambda_k$. Monotonicity on the two sides of
the gap gives
\[
 A_k^{3-k}\le(g\lambda_{k+1})^{k(3-k)},\qquad
 (\lambda_{k+1}\cdots\lambda_3)^k
 \ge\lambda_{k+1}^{k(3-k)}.
\]
Dividing these inequalities and using $k(3-k)=2$ yields
\begin{equation}\label{eq:partialminima}
 A_k\le P_\lambda^{k/3}g^{k(3-k)/3}
 \le8^{k/3}Q^{-\delta/3}.
\end{equation}

\emph{The determinant calculation.}
Let $F_1,\ldots,F_{R_0}$ be the distinct forms occurring in $\LL$.
They include the three coordinate forms, so $R_0\ge3$. Define the
fixed quantity
\[
 \mathfrak H=\prod_{v\in M_K}
 \max_{1\le i_1<i_2<i_3\le R_0}
       \|\det(F_{i_1},F_{i_2},F_{i_3})\|_v\ge1.
\]
This is the form height in \cite[Lemma~10.2]{EF13}.
Choose $0<\eta<1$ with $(1+\eta)\lambda_k<\lambda_{k+1}$ and
independent $g_1,\ldots,g_k\in T_k$ satisfying
$H_{\LL,c,Q}(g_j)\le(1+\eta)\lambda_j$. Such vectors can be chosen
successively from the definition of the infima. Work in a finite
extension containing their coordinates.

At each place $v$, select a restricted independent $k$-tuple
$L_{i_1(v),v},\ldots,L_{i_k(v),v}$ that realizes the minimum in the
definition of $w_c(T_k)$. For $w\mid v$ put
\[
 \theta_w=\det(L_{i_l(v),w}(g_j))_{l,j=1}^k,\qquad
 H_{jw}=\max_i\|L_{iw}(g_j)\|_wQ^{-c_{iw}}.
\]
These determinants are nonzero. Let $s(w)$ be the local degree divided
by $[E:\Q]$ at infinite places and zero at finite places; then
$\sum_ws(w)=1$. The determinant expansion gives
\[
 \|\theta_w\|_w\le(k!)^{s(w)}
 \prod_{j=1}^kH_{jw}\,Q^{\sum_{l=1}^kc_{i_l(v),w}}.
\]
Since $\sum_{w\mid v}d(w\mid v)=1$, multiplication gives
\begin{align}
 \prod_w\|\theta_w\|_w
 &\le k!\prod_{j=1}^kH_{\LL,c,Q}(g_j)Q^{w_c(T_k)}\notag\\
 &\le k!2^kA_k\le32Q^{-\delta/3}.
 \label{eq:detupper}
\end{align}
The middle inequality is where semistability is used.

There are at most $\binom{R_0}k$ distinct nonzero determinant values
obtained from all restricted $k$-tuples of the $F_i$. Set
$D_k=\binom{R_0}k-1$ and $C_{\rm form}=\sqrt3\,\mathfrak H$.
The minimum-determinant estimate \cite[Lemma~10.2, equation~(10.2)]{EF13}
states that
\begin{equation}\label{eq:detlower}
 \prod_w\min_{\theta\ne0}\|\theta\|_w
 \ge(C_{\rm form}H_2(T_k))^{-D_k}.
\end{equation}
In particular the left side is a lower bound for the product in
\eqref{eq:detupper}. Combining the two estimates gives
\[
 H_2(T_k)\ge C_{\rm form}^{-1}32^{-1/D_k}
                 Q^{\delta/(3D_k)}.
\]
With $D_* =\max(D_1,D_2)$, the condition
$Q^{\delta/6}\ge32C_{\rm form}^{D_*}$ absorbs the fixed factor and gives
\begin{equation}\label{eq:infimumheight}
 H_2(T_k)\ge Q^{\delta/(6D_k)}\ge Q^{c_0\delta},
 \qquad c_0=\frac1{12R_0^3}.
\end{equation}

\emph{Rescaling and exterior powers.}
Choose a full basis near the successive infima, with its first $k$
vectors spanning $T_k$. The rescaling and Davenport construction in
\cite[Lemmas~11.1--11.3]{EF13} gives, in a finite extension, a
permutation $\pi$ and a basis $h_1,h_2,h_3$ preserving these spans, with
\begin{align}
 \|L_{iw}(h_j)\|_w
 &\le3^{-s(w)}Q^{c_{iw}} &&(w\nmid v_0),\label{eq:davenportordinary}\\
 \|L_{\pi(i),w}(h_j)\|_w
 &\le(3^9\min(\lambda_i,\lambda_j))^{d(w\mid v_0)}
       &&(w\mid v_0).\label{eq:davenportspecial}
\end{align}
For this use of Davenport's lemma only the local rescaling and the
near-infimum basis are needed. Its proof chooses an arbitrarily small
positive error satisfying the gap conditions in its equation~(11.1).
It does not use the large cutoff in the final interval theorem.

Put $p=3-k$, and list the $p$-element subsets $I_1,I_2,I_3$ of
$\{1,2,3\}$ in lexicographic order. Define
\[
 \wt L_{lv}=\bigwedge_{i\in I_l}L_{iv},\qquad
 a_{lv}=\sum_{i\in I_l}c_{iv},\qquad
 \wt h_l=\bigwedge_{i\in I_l}h_i,
\]
and retain $z_1=\wt h_1,z_2=\wt h_2$. For $k=1$ these are
$h_1\wedge h_2,h_1\wedge h_3$; for $k=2$ they are $h_1,h_2$.
They are independent. Their span $T$ satisfies
$H_2(T)=H_2(T_k)$ by \cite[Lemma~6.1]{EF13}, proving
\eqref{eq:packetheight}.

Away from $v_0$, determinant expansion in
\eqref{eq:davenportordinary} gives the factor
$((p!)3^{-p})^{s(w)}\le1$, proving the ordinary bound in
\eqref{eq:packetlocal}. Each original weight lies in
$[-2\gamma_v,\gamma_v]$. When $p=2$, its pair sum is the negative
of the omitted weight. Thus $|a_{lv}|\le2\gamma_v$ in both cases;
also $\sum_la_{lv}=0$ and $a_{lv}=0$ outside $S$.

\emph{The exponents at the auxiliary place.}
Let $\nu_l=\prod_{i\in I_l}\lambda_i$, in the same order. From
\eqref{eq:davenportspecial} and the exterior determinant expansion,
the bounds for the retained vectors are, up to the permutation of rows,
\[
 (Q^{b_1},Q^{b_2},Q^{b_3})
 =(3^{27}\nu_1,3^{27}\nu_2,3^{27}\nu_2).
\]
The constant $3^{27}$ is a common upper bound for both exterior
degrees. We check the product separately:
\[
 \begin{array}{ll}
 k=1:&(\nu_1,\nu_2,\nu_3)
       =(\lambda_1\lambda_2,\lambda_1\lambda_3,\lambda_2\lambda_3),
       \quad\nu_1\nu_2^2=P_\lambda^2\lambda_1/\lambda_2,\\[2pt]
 k=2:&(\nu_1,\nu_2,\nu_3)=(\lambda_1,\lambda_2,\lambda_3),
       \quad\nu_1\nu_2^2=P_\lambda\lambda_2/\lambda_3.
 \end{array}
\]
Hence \eqref{eq:minima} and \eqref{eq:minimagap} give
$Q^{\sum_lb_l}\le64\cdot3^{81}Q^{-\delta/2}$.
If $Q^{\delta/6}\ge64\cdot3^{81}$, then $\sum_lb_l\le-\delta/3$.

For $p=1,2$ set $M_p=\max(C_0^p,8C_0^{3-p},1)$.
Every product of $p$ distinct infima satisfies
\[
 M_p^{-1}Q^{-p}\le\nu_l\le M_pQ^{3-p}.
\]
The lower bound uses the lower bound for each infimum in
\eqref{eq:minima}; the upper bound uses their total product and the
lower bound for the complementary infima. If
$Q\ge3^{27}\max(M_1,M_2)$, it follows that
$Q^{-3}\le3^{27}\nu_l\le Q^3$. This proves $|b_l|\le3$.

For completeness the conditions used above are exactly
\begin{equation}\label{eq:pointwisecutoffs}
 \begin{gathered}
 Q^\delta\ge3^{3/2},\qquad
 Q^{\delta/6}\ge32C_{\rm form}^{D_*},\\
 Q^{\delta/6}\ge64\cdot3^{81},\qquad
 Q\ge3^{27}\max(M_1,M_2).
 \end{gathered}
\end{equation}
They hold for $Q\ge\exp(A_0/\delta)$, $0<\delta\le1$, if we take
$A_0$ larger than each of
\[
 \tfrac32\log3,\quad6\log(32C_{\rm form}^{D_*}),\quad
 6\log(64\cdot3^{81}),\quad\log(3^{27}\max(M_1,M_2)).
\]
This completes the pointwise argument with constants depending only
on the fixed forms and field.
\end{proof}

\subsection{One auxiliary polynomial for the changing weights}
Suppose $m$ selected parameters have the same $k$ in
Lemma~\ref{lem:packet}, and label them $h=1,\ldots,m$.
Their exterior forms are identical. Write $a_{hlv},b_{hl}$ for their
exponents and put
\begin{equation}\label{eq:Gamma}
 \Gamma_v=\max_{1\le h\le m}\gamma_{hv},\qquad
 \gamma_{hv}=\max_i c^{(h)}_{iv}.
\end{equation}
Each $\gamma_{hv}$ is nonnegative and at most $1$, by admissibility.
Consequently
\begin{equation}\label{eq:Gammasum}
 \sum_{v\in S}\Gamma_v\le s.
\end{equation}
We use this bound rather than $1$: the maximum over $h$ is taken
separately at each place, so it cannot be moved outside the sum.
Ordinary vanishing conditions are needed only where $\Gamma_v>0$.

For positive block degrees $r_1,\ldots,r_m$, let $D_r=\sum_hr_h$ and
\[
 \mathcal U(r)=\{(j_{hl})\in\Z_{\ge0}^{mN}:\sum_lj_{hl}=r_h
                    \text{ for each }h\},\qquad
 V=|\mathcal U(r)|=\prod_h\binom{r_h+N-1}{N-1}.
\]
The elementary bound $\binom{r_h+N-1}{N-1}\le(eN)^{r_h}$ gives
\begin{equation}\label{eq:monomialcount}
 V\le(eN)^{D_r}.
\end{equation}

\begin{proposition}[Simultaneous coefficient vanishing]\label{prop:auxiliary}
Let $0<\alpha<1$ and suppose
\begin{equation}\label{eq:mcondition}
 m>2\alpha^{-2}\log(2s+2).
\end{equation}
There is a nonzero polynomial $P$ over $K$, homogeneous of degree
$r_h$ in the block $X_h\in K^N$, with the following properties.
In its expansion
$P=\sum_jd^{(v)}_j\prod_{h,l}\wt L_{lv}(X_h)^{j_{hl}}$, we have
\begin{align}
 d^{(v)}_j&=0 &&\text{if }\sum_{h,l}j_{hl}a_{hlv}/r_h\ge3m\alpha\Gamma_v,
       \quad v\in S,\ \Gamma_v>0,\label{eq:vanishordinary}\\
 d^{(v_0)}_j&=0 &&\text{if }\sum_{h,l}j_{hl}b_{hl}/r_h
                          \ge-m\delta/9+3m\alpha,\label{eq:vanishspecial}\\
 H_2(P)&\le C_K A^{D_r}.&&\label{eq:polyheight}
\end{align}
Here $C_K\ge1$ and $A\ge2$ are fixed, independent of $m,r_h$,
the weights, and the parameters.
\end{proposition}
\begin{proof}
\emph{Counting the equations.}
Choose $j$ uniformly from $\mathcal U(r)$. The blocks $j_h$ are
independent, and symmetry gives $\mathbb E(j_{hl})=r_h/N$.
For an ordinary place with $\Gamma_v>0$, the variables
\[
 X_{hv}=r_h^{-1}\sum_lj_{hl}a_{hlv}
\]
have mean zero and lie in $[-3\Gamma_v,3\Gamma_v]$.
Hoeffding's inequality \cite[Lemma~13.2]{EF13} therefore gives
\[
 \Pr\left(\sum_hX_{hv}\ge3m\alpha\Gamma_v\right)
 \le\exp\left(-\frac{2(3m\alpha\Gamma_v)^2}
                         {m(6\Gamma_v)^2}\right)
 =e^{-m\alpha^2/2}.
\]
At $v_0$ use $X_{h0}=r_h^{-1}\sum_lj_{hl}b_{hl}$.
Its mean is $N^{-1}\sum_lb_{hl}\le-\delta/9$, and its range is
$[-3,3]$. The event
$\sum_hX_{h0}\ge-m\delta/9+3m\alpha$ thus has probability at most
$e^{-m\alpha^2/2}$ as well.

Each coefficient required to vanish gives one homogeneous linear
equation in the $V$ coefficients of $P$ in the original variables.
Counting the equations separately in each coordinate system, their
total number $U$ satisfies
\begin{equation}\label{eq:equationcount}
 U\le(s+1)Ve^{-m\alpha^2/2}<V/2.
\end{equation}
This counts equations, including any redundancies, rather than a
union of forbidden monomial indices in different bases. Linear
dependencies can only reduce the number of constraints.

\emph{Uniform heights of the rows.}
Write $Y_h=\mathscr A_vX_h$, where $\mathscr A_v$ is the matrix of
the fixed exterior forms at $v$. The coefficient equations use the
substitution $X_h=\mathscr A_v^{-1}Y_h$. For a place $u$ of $K$ put
\[
 C_u=\max\bigl(1,\|\text{entries of }\mathscr A_v^{-1}\|_u:
                         v\in S\cup\{v_0\}\bigr),\qquad C=\prod_uC_u.
\]
The matrices are invertible and fixed, so $C$ is finite and fixed.
The special-place matrix is the identity. Here $s(u)$ is the local
degree divided by $[K:\Q]$ at an infinite place of $K$, and is zero
at a finite place, so $\sum_us(u)=1$. The absolute value of each
entry of a coefficient row is at most
$N^{D_rs(u)}C_u^{D_r}$, by the multinomial expansion.
Since the row has $V$ coordinates, its Euclidean
height is at most
\[
 V^{1/2}(NC)^{D_r}\le A_1^{D_r},\qquad
 A_1=(eN)^{1/2}NC.
\]
Thus changing weights changes which rows are selected, but does not
change their height bound.

If $U=0$, a coordinate monomial already suffices. Otherwise the
number-field Siegel lemma \cite[Lemma~13.1]{EF13} supplies a nonzero
coefficient vector with
\begin{align*}
 H_2(P)&\le C_KV^{1/2}
          \left(\prod_{u=1}^UH_2(\text{row}_u)\right)^{1/(V-U)}\\
 &\le C_K(eN)^{D_r/2}A_1^{D_rU/(V-U)}
 \le C_K\bigl((eN)^{1/2}A_1\bigr)^{D_r}.
\end{align*}
Taking $A\ge\max(2,(eN)^{1/2}A_1)$ proves \eqref{eq:polyheight}.
There are only two possible exterior degrees, so a common choice of
these constants works for both.
\end{proof}

\subsection{Divided derivatives and their local coefficients}
We differentiate in the original variables. For a multi-index $i$,
put
\[
 P_i=\left(\prod_{h,l}i_{hl}!\right)^{-1}\partial^iP,
 \qquad |i|_r=\sum_{h,l}i_{hl}/r_h.
\]
The following lemma makes the loss in the vanishing conditions and
the uniformity of the coefficient heights explicit.

\begin{lemma}[Derivative estimates]\label{lem:derivatives}
If $|i|_r\le2m\alpha$, the divided derivative has expansions
\begin{align}
 P_i&=\sum_jd^{(v)}_{i,j}\prod_{h,l}\wt L_{lv}(X_h)^{j_{hl}},
       \label{eq:derivativeexp}\\
 d^{(v)}_{i,j}&=0\quad\text{if }\sum_{h,l}j_{hl}a_{hlv}/r_h
                   >12m\alpha\Gamma_v,\quad v\in S,
                   \label{eq:derivativeordinary}\\
 d^{(v_0)}_{i,j}&=0\quad\text{if }\sum_{h,l}j_{hl}b_{hl}/r_h
                  >-m\delta/9+12m\alpha,\label{eq:derivativespecial}\\
 \prod_v\max_j\|d^{(v)}_{i,j}\|_v&\le C_K A_2^{D_r},
       \label{eq:derivativeheight}
\end{align}
where $A_2\ge2$ is fixed. Every monomial in a nonzero derivative
has block degree $r_h-\sum_li_{hl}$.
\end{lemma}
\begin{proof}
In the coordinates $Y_h=\mathscr A_vX_h$, the chain rule reads
\[
 \frac{\partial}{\partial X_{ha}}
 =\sum_l(\mathscr A_v)_{la}\frac{\partial}{\partial Y_{hl}}.
\]
It follows that a coefficient with index $j$ in the derivative
is a linear combination of original $Y$-coefficients with indices
$j+q$, where $q_{hl}\ge0$ and
$\sum_lq_{hl}=\sum_li_{hl}$ in each block. The ordinary weighted
shift satisfies
\[
 \left|\sum_{h,l}q_{hl}a_{hlv}/r_h\right|
 \le2\Gamma_v|i|_r\le4m\alpha\Gamma_v
 \le6m\alpha\Gamma_v.
\]
At the special place the same shift is at most
$3|i|_r\le6m\alpha$ in absolute value.
If either strict inequality in \eqref{eq:derivativeordinary} or
\eqref{eq:derivativespecial} holds, every contributing original
coefficient therefore meets its vanishing condition in
Proposition~\ref{prop:auxiliary}. This proves the assertions of
vanishing. When $\Gamma_v=0$, all ordinary weights at $v$ are zero,
so the strict inequality there cannot occur.

For the height estimate let $a_P$ be the original coefficient vector
of $P$. In the original variables, a divided derivative multiplies
each coefficient by a product of binomial coefficients, bounded by
$2^{D_r}$ at infinite places and by $1$ at finite places. Passing
to the $Y$-coordinates uses the fixed inverse matrix. At any place $u$
and for any one of the relevant coordinate systems this gives
\[
 \max_j\|d^{(v)}_{i,j}\|_u
 \le V^{s(u)}(2N)^{D_rs(u)}C_u^{D_r}
                       \max_j\|(a_P)_j\|_u.
\]
For systems outside $S\cup\{v_0\}$ the matrix is the identity, and
the same estimate remains valid. Multiply this bound using the
system assigned to each place. Since $\sum_us(u)=1$ and the sup
height of $a_P$ is at most $H_2(P)$, \eqref{eq:monomialcount} gives
\[
 \prod_v\max_j\|d^{(v)}_{i,j}\|_v
 \le(2eN^2C)^{D_r}H_2(P)
 \le C_K\bigl(2eN^2CA\bigr)^{D_r}.
\]
This proves \eqref{eq:derivativeheight} with a fixed $A_2$.
\end{proof}

\subsection{Degree selection and the nonvanishing condition}
\begin{proof}[Proof of Theorem~\ref{thm:sparse}]
Choose
\begin{equation}\label{eq:parameters}
 \alpha=\frac{\delta}{1000(s+1)},\qquad
 m=\left\lfloor2\alpha^{-2}\log(2s+2)\right\rfloor+1,
 \qquad \Omega=\frac{4m^2}{\alpha}.
\end{equation}
We specify $C_*(\delta)$ below; in particular it will be at least
$\exp(A_0/\delta)$. Lemma~\ref{lem:packet} therefore applies to
every parameter. There are two possible values of $k$, so a sequence
of at least $2m-1$ parameters has $m$ with a common $k$.
Separation is preserved on passing to this subsequence. Relabel its
parameters $Q_1<\cdots<Q_m$ and its exterior data by $h$.
All vectors can be placed in a common finite extension of $K$;
the normalized local bounds are unchanged by further extension.

Choose an arbitrarily large positive integer $r_1$ such that
$r_1\log Q_1>\alpha^{-1}\log Q_m$ and $C_K\le2^{r_1}$.
For $h=1,\ldots,m$ put
\[
 L_h=\frac{r_1\log Q_1}{\log Q_h},\qquad
 r_h=\lceil L_h\rceil,\qquad B=Q_1^{r_1}.
\]
This agrees with the chosen integer $r_1$ at $h=1$.
Since $L_h>1/\alpha$, we have
$L_h\le r_h<L_h+1<(1+\alpha)L_h$. Also $L_h\le r_1$, hence
$r_h\le r_1$. These facts give
\begin{equation}\label{eq:degrees}
 B\le Q_h^{r_h}<B^{1+\alpha},\qquad r_h\le r_1,\qquad
 \frac{r_h}{r_{h+1}}>
 \frac{\log Q_{h+1}}{(1+\alpha)\log Q_h}
 >\frac{\Omega}{1+\alpha}\ge\frac{2m^2}{\alpha}.
\end{equation}
The last ratio is in the decreasing-degree direction required by
\cite[Proposition~12.1, equation~(12.3)]{EF13}.

Construct $P$ by Proposition~\ref{prop:auxiliary}. In each plane $T_h$
take the grid
\[
 \mathcal G_h=\{u_1z_{h1}+u_2z_{h2}:u_j\in\Z,
                    |u_j|\le\lceil N/\alpha\rceil\},\qquad N=3.
\]
It contains a grid of the size $N/\alpha$ required by that
nonvanishing proposition. Set
\[
 E_m=(N-1)(3m^2/\alpha)^m,\qquad A_3=2eA.
\]
By \eqref{eq:polyheight} and $C_K\le2^{r_1}\le2^{D_r}$,
we have $e^{D_r}H_2(P)\le A_3^{D_r}$.
The remaining hypothesis of the proposition is
\begin{equation}\label{eq:nonvanishheight}
 H_2(T_h)^{r_h}\ge(e^{D_r}H_2(P))^{E_m}
 \quad(h=1,\ldots,m).
\end{equation}
Its two sides have the bounds
\[
 H_2(T_h)^{r_h}\ge Q_h^{c_0\delta r_h}\ge B^{c_0\delta},
 \qquad
 (e^{D_r}H_2(P))^{E_m}\le A_3^{D_rE_m}\le A_3^{mr_1E_m}.
\]
It is therefore sufficient that
\begin{equation}\label{eq:cutoffone}
 c_0\delta\log Q_1\ge mE_m\log A_3.
\end{equation}
Under this condition \cite[Proposition~12.1]{EF13} supplies
$x_h\in\mathcal G_h\setminus\{0\}$ and a multi-index with
$|i|_r\le2m\alpha$ for which
\[
 P_i(x_1,\ldots,x_m)\ne0.
\]
We now estimate this nonzero algebraic number at every place.

\emph{The grid factor.}
At an infinite place $w$, the triangle inequality introduces at most
the sum of the absolute values of the two grid coefficients.
In general, with $N-1$ coefficients, this sum is at most
\[
 (N-1)\lceil N/\alpha\rceil
 \le(N-1)(N/\alpha+1)\le N^2/\alpha,
\]
because $\alpha\le1$. At finite places the coefficients are integers
and the ultrametric inequality introduces no factor. Thus
\eqref{eq:packetlocal} gives
\begin{equation}\label{eq:gridfactor}
 \|\wt L_{lw}(x_h)\|_w\le
 (N^2/\alpha)^{s(w)}Q_h^{d(w\mid v)a_{hlv}}
 \quad(w\mid v\ne v_0).
\end{equation}
At $w\mid v_0$ use $b_{hl}$ instead, with no grid factor since $v_0$
is finite. The factor $N^2/\alpha$ will be included in the final bound.

\emph{Ordinary-place exponents.}
Put $\kappa_h=r_h\log Q_h/\log B$, so $1\le\kappa_h<1+\alpha$.
For a surviving monomial in the derivative at $w\mid v\ne v_0$, set
\[
 A_{hw}=\frac1{r_h}\sum_lj_{hl}d(w\mid v)a_{hlv},\qquad
 \Gamma_w=d(w\mid v)\Gamma_v.
\]
The block degrees are at most $r_h$, so $|A_{hw}|\le3\Gamma_w$.
The parameter contribution to this monomial is exactly
\[
 \prod_hQ_h^{r_hA_{hw}}=B^{\sum_h\kappa_hA_{hw}}.
\]
Using the derivative vanishing condition and then the ceiling error,
\[
 \sum_hA_{hw}\le12m\alpha\Gamma_w,\qquad
 \left|\sum_h(\kappa_h-1)A_{hw}\right|
 \le3m\alpha\Gamma_w.
\]
The ordinary-place exponent is therefore at most
$15m\alpha\Gamma_w$. This also holds outside $S$, where all the
ordinary exponents vanish.

\emph{The exponent at $v_0$.}
Here define $A_{hw}=d(w\mid v_0)\sum_lj_{hl}b_{hl}/r_h$.
We have $|A_{hw}|\le3d(w\mid v_0)$, while
\eqref{eq:derivativespecial} gives
\[
 \sum_hA_{hw}\le d(w\mid v_0)(-m\delta/9+12m\alpha).
\]
The ceiling error is at most $3m\alpha d(w\mid v_0)$.
The special-place exponent is thus at most
$d(w\mid v_0)(-m\delta/9+15m\alpha)$.

\emph{Multiplication over all places.}
A derivative has at most $V\le(eN)^{D_r}$ monomials. Locally the sum
introduces the factor $V^{s(w)}$, and the grid contributes at most
$(N^2/\alpha)^{D_rs(w)}$. Multiply these with
\eqref{eq:derivativeheight}. The identities
\[
 \sum_ws(w)=1,\qquad \sum_{w\mid v}d(w\mid v)=1,
 \qquad\sum_{v\in S}\Gamma_v\le s
\]
give
\begin{align}
 \prod_w\|P_i(x_1,\ldots,x_m)\|_w
 &\le C_K(A_2eN^3/\alpha)^{D_r}
       B^{m[-\delta/9+15(s+1)\alpha]}\notag\\
 &\le G(\alpha)^{D_r}(B^m)^{-\delta/12},
       \qquad G(\alpha)=2A_2eN^3/\alpha.
       \label{eq:product}
\end{align}
Here $C_K\le2^{D_r}$, and our choice of $\alpha$ ensures
\[
 -\delta/9+15(s+1)\alpha
 =-\delta/9+15\delta/1000\le-\delta/12.
\]
Since $D_r\le mr_1$ and $B=Q_1^{r_1}$, the logarithm of the right
side of \eqref{eq:product} is at most
\[
 mr_1\bigl(\log G(\alpha)-(\delta/12)\log Q_1\bigr).
\]
It is strictly negative if
\begin{equation}\label{eq:cutofftwo}
 \log Q_1\ge24\delta^{-1}\log G(\alpha).
\end{equation}
This contradicts the product formula for the nonzero value supplied
by the nonvanishing proposition.

\emph{The uniform cutoff.}
To meet the pointwise cutoff, \eqref{eq:cutoffone}, and
\eqref{eq:cutofftwo} simultaneously, take
\begin{equation}\label{eq:explicitcutoff}
 C_*(\delta)=\exp\left(
 \frac{A_0+1}{\delta}
 +\frac{2mE_m\log A_3}{c_0\delta}
 +\frac{24\log G(\alpha)}{\delta}\right).
\end{equation}
The dependence on $\delta$ can now be checked without suppressing
the nonvanishing exponent. With the fixed support, we have
$\alpha\asymp\delta$, $m\asymp\delta^{-2}$, and consequently
$\Omega=4m^2/\alpha\ll\delta^{-5}$. Moreover
\begin{equation}\label{eq:nonvanishingexponent}
 \log E_m=\log(N-1)+m\log(3m^2/\alpha)
 \ll\delta^{-2}\log(2/\delta).
\end{equation}
Taking the logarithm of the sum defining $\log C_*$, and bounding
it by the largest summand up to a fixed factor, gives
\[
 \log\log C_*
 \ll1+\log m+\log E_m+\log(1/\delta)
          +\log(1+\log G(\alpha))
 \ll\delta^{-2}\log(2/\delta).
\]
All fixed factors are absorbed into constants depending only on the
fixed data. A sequence of length $2m-1$ has been ruled out, so its
length is at most $2m-2=O(\delta^{-2})$. This proves the theorem.
\end{proof}

\begin{corollary}\label{cor:count}
Under the same hypotheses, suppose $Q_{h+1}>Q_h^2$.
The number of parameters at or above $C_*(\delta)$ is
$O(\delta^{-2}\log(2/\delta))$.
\end{corollary}
\begin{proof}
Discard any parameters below $C_*(\delta)$ and set
$q=\lceil\log_2\Omega\rceil\ll\log(2/\delta)$.
Partition the remaining indices into their $q$ residue classes.
Successive parameters in any one class are $q$ steps apart, and
iteration of the strict doubling inequality gives
$Q_{h+q}>Q_h^{2^q}\ge Q_h^\Omega$.
Theorem~\ref{thm:sparse} bounds each class by $O(\delta^{-2})$;
summing the $q$ bounds proves the assertion.
\end{proof}

The logarithm in this corollary comes from the separation partition.
The logarithm in \eqref{eq:cutofforder} comes from
\eqref{eq:nonvanishingexponent}. These are distinct conditions, and
both will be retained in the dyadic counting argument below.

\section{Dyadic scales and fixed-quality layers}\label{sec:layers}
Fix the forms and support from Section~\ref{sec:height}, and choose
a finite place $v_0$ above a rational prime not dividing $b$.
The admissibility requirements in Section~\ref{sec:moving} follow
from \eqref{eq:norm} and Proposition~\ref{prop:stable}.

\begin{lemma}[Counting all prefix scales]\label{lem:layers}
There are constants $K_0,K_1>0$, depending only on $\xi,b$, with the
following property. Let $B\ge2$ and let $I$ be a finite set of integers.
For each $i\in I$, suppose there is an approximation satisfying
\eqref{eq:repeat}--\eqref{eq:approx}, with $H=2^i$, such that
\[
 t_i\ge T,\qquad v_i\ge2^i/B.
\]
If
\begin{equation}\label{eq:layerthreshold}
 \log T\ge K_0B^2\log(2B),
\end{equation}
then
\begin{equation}\label{eq:layercount}
 |I|\le K_1B^2\log(2B).
\end{equation}
\end{lemma}
\begin{proof}
Put $\tau_i=\lfloor\log_2t_i\rfloor$ and
$\mu_\tau=\#\{i\in I:\tau_i=\tau\}$.
For each nonempty class retain its largest index $i_\tau$ and that
index's approximation. Since $t_i\le2^i$, every member satisfies
$i\ge\tau$. The indices are distinct, so
$i_\tau\ge\tau+\mu_\tau-1$. As $t_{i_\tau}<2^{\tau+1}$, this yields
\begin{equation}\label{eq:multiplicityquality}
 \frac{v_{i_\tau}}{t_{i_\tau}}
 >\frac{2^{i_\tau-\tau-1}}B
 \ge\frac{2^{\mu_\tau-2}}B.
\end{equation}

For each integer $j\ge1$ set
\[
 \mathcal T_j=\{\tau:\mu_\tau\ge j\},\quad N_j=|\mathcal T_j|,
 \qquad \eps_j=\frac{2^{j-3}}B,\quad
 \delta_j=\frac{\eps_j}{3+\eps_j}.
\]
Only levels with $N_j>0$ need consideration. At such a level,
\eqref{eq:multiplicityquality} and $v_i/t_i\le1$ imply
$2^{j-2}/B<1$, hence $0<\eps_j<1/2$.
Every retained approximation in that level has $v_i/t_i>2\eps_j$,
so \eqref{eq:small} applies with the common $\eps_j$ and $\delta_j$.

Partition $\mathcal T_j$ according to the parity of $\tau$.
In either part consecutive values differ by at least two. Thus their
retained parameters satisfy
$t_{\mathrm{next}}\ge2^{\tau+2}>2t_{\mathrm{previous}}$,
and consequently $Q_{\mathrm{next}}>Q_{\mathrm{previous}}^2$.

We verify the cutoff uniformly over all nonempty levels.
We have $\delta_j\ge\delta_1=1/(12B+1)$.
The function $\delta^{-2}\log(2/\delta)$ decreases on $(0,1]$;
therefore \eqref{eq:cutofforder} gives
\[
 \log\log C_*(\delta_j)\ll B^2\log(2B).
\]
On the other hand, $Q=b^{D_{\eps_j}t}$ and $D_{\eps_j}\ge1$ give
$\log\log Q\ge\log T+\log\log b$.
Increasing $K_0$ in \eqref{eq:layerthreshold} puts every retained
point above its level's cutoff. Monotonicity of $C_*$ itself is not needed.

Apply Corollary~\ref{cor:count} to both parity classes. Since
$\delta_j\ge2\eps_j/7$, it follows that
\begin{equation}\label{eq:tailcount}
 N_j\ll\delta_j^{-2}\log(2/\delta_j)
 \ll B^2 4^{-j}\log(2B).
\end{equation}
The constants are uniform in $B,j$ and the approximations.
Summation over multiplicity levels now gives
\[
 |I|=\sum_\tau\mu_\tau=\sum_{j\ge1}N_j
 \ll B^2\log(2B)\sum_{j\ge1}4^{-j}
 \ll B^2\log(2B),
\]
as required.
\end{proof}

\begin{remark}
Each application of Corollary~\ref{cor:count} uses a fixed quality
$\delta_j$. Changing the layer changes this common quality and the
associated normalized heights. The uniform bound for the cutoffs
makes this legitimate; no theorem for a chain with varying
smallness exponents $\delta_h$ is used.
\end{remark}

\section{Window bounds and the endpoint}\label{sec:completion}
\begin{proof}[Proof of Theorem~\ref{thm:main}]
Fix $\theta>1$, let $Y=X^\theta$, and put
$F(x)=\sqrt{\log x/\log\log x}$.
Suppose that, throughout the integer lengths in $[X,Y]$,
\begin{equation}\label{eq:windowLC}
 R(k)\le CkF(k),
\end{equation}
where $C>0$ is a fixed constant to be chosen sufficiently small.
Let $F_Y=F(Y)$, $B=24CF_Y$, and
\[
 I=\{i\in\Z:8CF_YX\le2^i\le Y\}.
\]
For sufficiently large $X$, $B\ge2$ and $CF_Y\ge1$.
Lemma~\ref{lem:windowwords} provides the approximations needed in
Lemma~\ref{lem:layers}, with $T=X/3$.
The window buffer ensures that every invoked value of $R(k)$
lies inside $[X,Y]$.

For fixed $C,\theta$, the number of scales satisfies
\begin{equation}\label{eq:windowscales}
 |I|=\frac{\log(Y/X)}{\log2}-O(\log F_Y+|\log C|+1)
 =\left(\frac{\theta-1}{\log2}+o(1)\right)\log X.
\end{equation}
Also
\begin{equation}\label{eq:windowbudget}
 B^2\log(2B)
 =(288C^2+o(1))\log Y
 =(288\theta C^2+o(1))\log X.
\end{equation}
Choose $C>0$ so small that
\[
 288\theta K_0C^2<\tfrac12,\qquad
 288\theta K_1C^2<\frac{\theta-1}{2\log2}.
\]
Then $\log(X/3)\ge K_0B^2\log(2B)$ for sufficiently large $X$,
and Lemma~\ref{lem:layers} applies.
Its bound \eqref{eq:layercount} contradicts
\eqref{eq:windowscales}--\eqref{eq:windowbudget}.
Hence some integer $L\in[X,X^\theta]$ has $R(L)>CLF(L)$.
By \eqref{eq:RvsP}, the same $L$ has
$p(L)\ge R(L)-L>(C/2)LF(L)$ once $X$ is large enough.
This proves the theorem with $c=C$.
\end{proof}

\begin{proof}[Proof of Corollary~\ref{cor:limsup}]
Take, for example, $\theta=3/2$ in Theorem~\ref{thm:main}.
The lengths provided tend to infinity with $X$, proving the positive
limsup. Multiplying the resulting lower bounds by
$(\log\log L)^{\gamma-1/2}$ proves the first assertion in
\eqref{eq:weightedlimsup}; multiplying by
$(\log L)^{1/2-\eta}/\sqrt{\log\log L}$ proves the second.
\end{proof}

\begin{proof}[Proof of Theorem~\ref{thm:nearwindow}]
For the given $F$, we have $F(X)\to\infty$ and
\begin{equation}\label{eq:subcritical}
 F(X)^2\log(2F(X))=o(\log X).
\end{equation}
Fix $C_0>0$ and write
\[
 W=A F(X)^2\log(2F(X)),\qquad Y=Xe^W.
\]
For fixed $A$, equation \eqref{eq:subcritical} gives $W=o(\log X)$.
Thus $\log Y/\log X\to1$. Substitution into the definition of $F$
then yields $F(Y)/F(X)\to1$.

Suppose $R(k)\le C_0kF(k)$ at all integer lengths in $[X,Y]$.
Apply Lemma~\ref{lem:windowwords} with $B=24C_0F(Y)$ and
$I=\{i:8C_0F(Y)X\le2^i\le Y\}$.
Since $W/\log F(X)\to\infty$, we have
\[
 |I|=\left(\frac A{\log2}+o(1)\right)F(X)^2\log(2F(X)),
 \qquad
 B^2\log(2B)=(576C_0^2+o(1))F(X)^2\log(2F(X)).
\]
The threshold \eqref{eq:layerthreshold}, with $T=X/3$, holds by
\eqref{eq:subcritical}. Choose $A>576K_1C_0^2\log2$.
For all sufficiently large $X$, \eqref{eq:layercount} is then a
contradiction. Thus some $L\in[X,Y]$ satisfies $R(L)>C_0LF(L)$.
To obtain the assertion for the prescribed $C$, apply this with
$C_0=2C$. For large $X$, \eqref{eq:RvsP} gives
$p(L)>CLF(L)$ at the same $L$, as well as $R(L)>CLF(L)$.
Finally \eqref{eq:subcritical} gives $Y=X^{1+o(1)}$.
\end{proof}

\section{Scope and the remaining losses}
The arithmetic proof uses finite common support to impose simultaneous
coefficient vanishing with $m=O(\delta^{-2})$, and an explicit check of
semistability for the full family of exact weights. In particular,
the moving-weight theorem is proved in Section~\ref{sec:moving}; it
does not follow by simply applying the published fixed-weight theorem
at each point. Nonreal conjugates and composite bases are handled by
the mass and prime-weight identities in Section~\ref{sec:height}.

The dyadic argument records multiplicity across prefix scales.
It does not give an arithmetic bound for all pairs of equal words
within one prefix. If $f_w$ counts the occurrences of a length-$k$
word in a prefix of length $H$, then Cauchy--Schwarz gives
\[
 \sum_w f_w^2\ge\frac{(H-k+1)^2}{p(k)}.
\]
The repetition lemma only needs one pair; multiplicities of arbitrary
pairs are not retained by the subsequent arithmetic construction.
Lemma~\ref{lem:layers} instead retains the multiplicity of prefix
scales through the quality gain \eqref{eq:multiplicityquality}.

The exponent $1/2$ in the present estimates comes from
\eqref{eq:mcondition}, together with $\alpha\asymp\delta$ in
\eqref{eq:parameters}. Two additional restrictions each contribute
one logarithm at the final scale: the nonvanishing threshold has
$\log\log C_*\ll\delta^{-2}\log(2/\delta)$, and passing from
$\Omega$-separation to ordinary doubling separation costs
$q\asymp\log(2/\delta)$ in Corollary~\ref{cor:count}.
These are separate requirements at the same order
$B^2\log(2B)$; their logarithmic costs are not multiplied in the
dyadic proof. Both requirements must be addressed to remove the
remaining $\sqrt{\log\log L}$ denominator by this method.
This identifies limitations of the argument given here, without
asserting an obstruction to all other methods.

The endpoint proof gives fixed-power windows with constants depending
on their exponent. The proof of shrinking relative logarithmic
windows uses \eqref{eq:subcritical}; it does not establish such windows
at $u=\gamma=1/2$. No conclusion of positive entropy, disjunctivity,
or normality follows from these block-complexity lower bounds.

\begin{thebibliography}{9}
\bibitem{AB07} B.~Adamczewski and Y.~Bugeaud,
On the complexity of algebraic numbers I. Expansions in integer bases,
\emph{Ann. of Math.} (2) \textbf{165} (2007), no.~2, 547--565.
\href{https://doi.org/10.4007/annals.2007.165.547}{doi:10.4007/annals.2007.165.547}.
\bibitem{BE08} Y.~Bugeaud and J.-H.~Evertse,
On two notions of complexity of algebraic numbers,
\emph{Acta Arith.} \textbf{133} (2008), no.~3, 221--250.
\href{https://doi.org/10.4064/aa133-3-3}{doi:10.4064/aa133-3-3}.
\bibitem{EF13} J.-H.~Evertse and R.~G.~Ferretti,
A further improvement of the quantitative Subspace Theorem,
\emph{Ann. of Math.} (2) \textbf{177} (2013), no.~2, 513--590.
\href{https://doi.org/10.4007/annals.2013.177.2.4}{doi:10.4007/annals.2013.177.2.4}.
\bibitem{BK19} Y.~Bugeaud and D.~H.~Kim,
A new complexity function, repetitions in Sturmian words, and
irrationality exponents of Sturmian numbers,
\emph{Trans. Amer. Math. Soc.} \textbf{371} (2019), no.~5, 3281--3308.
\href{https://doi.org/10.1090/tran/7378}{doi:10.1090/tran/7378}.
\end{thebibliography}
\end{document}
